\begin{abstract}
Orbital missions have mapped the Moon with many instruments, from multispectral images taken under changing sunlight to static maps of topography, roughness and composition. The products are co-registered but differ in resolution, units and physics, and labels are scarce. We present \method, a self-supervised model that learns from this stack by predicting, in representation space, what one instrument would record at a 60\,km footprint given another instrument's view of it. Per-product stems give every product the same ground cell per token, from 60\,m to 1\,km per pixel, and a cross-attention predictor, told the target product and its sun geometry, predicts the token grid of a moving-average encoder. Training this way has a failure mode that the loss hides. Several runs converged to a position code, in which tokens depend on their place in the grid and hardly on the terrain, while the loss fell tenfold and within-image rank and VICReg penalties improved. We explain why these signals are blind, show that a covariance penalty and a variance hinge on predictions push toward the shortcut, and propose two cross-sample monitors that flagged all seven failures we analyzed. Frozen \method features detect craters on par with a masked-token lunar foundation model and gain up to 0.23 mAP$_{50}$ from a second input product. Probes recover relief and albedo from single images, and predicted shadows follow the sun the predictor is given (AUC 0.89 with the true sun, 0.59 with another).
\end{abstract}
